\newcommand{\eqname}[1]{\text{(#1)}\qquad}
\newcommand{\eqnote}[1]{\tag*{(#1)}}

\newcommand{\composition}{\circ}
\newcommand{\define}{=\vcentcolon}
\newcommand{\defined}{\vcentcolon=}
\newcommand{\isomorphic}{\cong}
\newcommand{\complexes}{\mathbb{C}}
\newcommand{\reals}{\mathbb{R}}
\newcommand{\integers}{\mathbb{Z}}
\newcommand{\naturals}{\mathbb{N}}
\newcommand{\continuous}{C}
\newcommand{\euclideans}{\mathbb{E}}
\newcommand{\Mat}[1][]{\ifthenelse{\equal{#1}{}}{\textnormal{Mat}}{\textnormal{Mat}(#1)}}
\newcommand{\SO}[1]{\textnormal{SO}({#1})}
\newcommand{\SE}[1]{\textnormal{SE}({#1})}
\renewcommand{\so}[1]{\mathfrak{so}(#1)}
\newcommand{\se}[1]{\mathfrak{se}(#1)}
\newcommand{\RSO}[1]{R_{\textnormal{\tiny{$\SO$}}}\hspace{-3pt}\left(#1\right)}
\newcommand{\RSE}[1]{R_{\textnormal{\tiny{$\SE$}}}\hspace{-3pt}\left(#1\right)}
\newcommand{\tangent}[1]{
    \ifthenelse{\equal{#1}{}}
    {{T}}
    {{T_{#1}}}
}
\newcommand{\dualtangent}[1]{
    \ifthenelse{\equal{#1}{}}
    {{T^*}}
    {{T_{#1}^*}}
}

\newcommand{\leftaction}[1]{L_{#1}}
\newcommand{\leftextension}[2][e]{\tangent{#1}\leftaction{#2}}
\newcommand{\dualleftextension}[2][e]{\dualtangent{#1}\leftaction{#2}}
\newcommand{\leftextensionfromvector}[1]{\Xi_{#1}}
\newcommand{\MaurerCartanform}[1]{\chi_{#1}}
% \newcommand{\TSE}[2]{
%     \ifthenelse{\equal{#1}{}}
%     {{T\SE{#2}}}
%     {{T_{#1}\SE{#2}}}
% }
% \newcommand{\TSE}[2]{
%     \tangent[#1]\SE{#2}
% }
% \newcommand{\dualTSE}[2][]{
%     \ifthenelse{\equal{#1}{}}
%     {{\dualtangent\SE{#2}}}
%     {{\dualtangent[#1]\SE{#2}}}
% }
\newcommand{\dualse}[1]{\mathfrak{se}^*(#1)}
\newcommand{\sebasis}{X}
\newcommand{\Adjoint}[1]{
    \ifthenelse{\equal{#1}{}}
    {\textnormal{Ad}}
    {\textnormal{Ad}_{#1}}
}

\newcommand{\adjoint}[1]{
    \ifthenelse{\equal{#1}{}}
    {\textnormal{ad}}
    {\textnormal{ad}_{#1}}
}
\newcommand{\coAdjoint}[1]{
    \ifthenelse{\equal{#1}{}}
    {\textnormal{Ad}^*}
    {\textnormal{Ad}^*_{#1}}
}

\newcommand{\coadjoint}[1]{
    \ifthenelse{\equal{#1}{}}
    {\textnormal{ad}^*}
    {\textnormal{ad}^*_{#1}}
}
\newcommand{\maptoliealgebra}[1]{\left(#1\right)_{\wedge}}
\newcommand{\maptovector}[1]{\left(#1\right)_{\vee}}
\newcommand{\maptodualliealgebra}[1]{\left(#1\right)^*_{\wedge}}

\newcommand{\perpendicular}{\bot}

\newcommand{\circconvolute}[1][\ast]{\raisebox{1pt}{\scriptsize{\textcircled{\tiny{$#1$}}}}}
\newcommand{\convolute}{*}

\newcommand{\dif}{\mathrm{d}}
\newcommand{\Dif}{\mathrm{D}}
\newcommand{\gradient}{\nabla}
\newcommand{\divergence}{\nabla\cdot}
\newcommand{\Hessian}{\mathbf{H}}

\newcommand{\zeros}{\mathbf{0}}
\newcommand{\ones}{\mathbf{1}}
\newcommand{\Set}[1]{\left\{#1\right\}}
\newcommand{\given}{\,\middle|\,}
\newcommand{\at}{|}
\renewcommand{\exp}[1]{\textnormal{ exp}\left(#1\right)}
\newcommand{\csch}[1]{\textnormal{ csch}\left(#1\right)}
\newcommand{\sech}[1]{\textnormal{ sech}\left(#1\right)}
\newcommand{\softmax}[1]{\textnormal{softmax}\left(#1\right)}

\newcommand{\sgn}[1]{\textnormal{sgn}\left(#1\right)}
\newcommand{\abs}[1]{\left|#1\right|}
\newcommand{\norm}[1]{\left\lVert#1\right\rVert}
\renewcommand{\det}[1]{\textnormal{det}\left(#1\right)}
\newcommand{\inverse}{{-1}}
\newcommand{\pseudoinverse}{+}
\newcommand{\conjugate}{*}
\newcommand{\transpose}{\intercal}
\newcommand{\conjugatetranspose}{\dagger}
\newcommand{\Trace}[1]{\textnormal{Tr}\left(#1\right)}
\newcommand{\diag}[1]{\textnormal{diag}\left\{#1\right\}}
\newcommand{\inner}[2]{\left\langle #1, #2 \right\rangle}

% Define the double angle brackets without using MnSymbol package
% The MnSymbol package causes weird symbol issues, e.g. \int
\makeatletter
\DeclareFontFamily{OMX}{MnSymbolE}{}
\DeclareSymbolFont{MnLargeSymbols}{OMX}{MnSymbolE}{m}{n}
\SetSymbolFont{MnLargeSymbols}{bold}{OMX}{MnSymbolE}{b}{n}
\DeclareFontShape{OMX}{MnSymbolE}{m}{n}{
    <-6>  MnSymbolE5
   <6-7>  MnSymbolE6
   <7-8>  MnSymbolE7
   <8-9>  MnSymbolE8
   <9-10> MnSymbolE9
  <10-12> MnSymbolE10
  <12->   MnSymbolE12
}{}
\DeclareFontShape{OMX}{MnSymbolE}{b}{n}{
    <-6>  MnSymbolE-Bold5
   <6-7>  MnSymbolE-Bold6
   <7-8>  MnSymbolE-Bold7
   <8-9>  MnSymbolE-Bold8
   <9-10> MnSymbolE-Bold9
  <10-12> MnSymbolE-Bold10
  <12->   MnSymbolE-Bold12
}{}

\let\llangle\@undefined
\let\rrangle\@undefined
\DeclareMathDelimiter{\llangle}{\mathopen}%
                     {MnLargeSymbols}{'164}{MnLargeSymbols}{'164}
\DeclareMathDelimiter{\rrangle}{\mathclose}%
                     {MnLargeSymbols}{'171}{MnLargeSymbols}{'171}
\makeatother
\newcommand{\pairing}[2]{\left\llangle #1, #2 \right\rrangle}
\newcommand{\liebracket}[2]{\left[ #1, \ #2 \right]}
\newcommand{\Frobenius}{\textnormal{F}}
\newcommand{\Frobeniusnorm}[1]{\norm{#1}_\Frobenius}

\newcommand{\identity}[1]{
    \ifthenelse{\equal{#1}{}}
    {\mathrm{I}}
    {\mathrm{I}_{#1}}
}

\newcommand{\union}{\bigcup}

\newcommand{\ProbabilitySimplex}{\Delta}
\newcommand{\ProbabilitySpace}[1]{\mathcal{P}(#1)}
\newcommand{\sigmaAlgebra}[1]{\sigma\left(#1\right)}
\newcommand{\indicator}[1]{\textbf{1}_{\left\{#1\right\}}}
\NewDocumentCommand{\Probability}{o m}{
  \mathsf{P}
  \IfValueT{#1}{_{#1}}
  \IfNoValueT{#1}{}
  \left(#2\right)
}
\NewDocumentCommand{\Expectation}{o m}{
  \mathsf{E}
  \IfValueT{#1}{_{#1}}
  \IfNoValueT{#1}{}
  \left(#2\right)
}
\NewDocumentCommand{\Variance}{o m}{
  \mathsf{Var}
  \IfValueT{#1}{_{#1}}
  \IfNoValueT{#1}{}
  \left(#2\right)
}
\DeclareMathOperator*{\logsumexp}{logsumexp}
\newcommand{\KLDivergence}[2]{\textnormal{KL}\left(#1 \| #2\right)}
\newcommand{\entropy}[1]{\mathsf{H}\left(#1\right)}

\newcommand{\Dim}{\textnormal{Dim}}
\renewcommand{\dim}{\mathsf{d}}

\DeclareMathOperator*{\argmax}{arg\,max}
\DeclareMathOperator*{\argmin}{arg\,min}

\newcommand{\subjectto}{\textnormal{s.t.}}

\newcommand{\Order}[1]{\mathcal{O}\left(#1\right)}

\newcommand{\mysetminusD}{\hbox{\tikz{\draw[line width=0.6pt,line cap=round] (3pt,0) -- (0,6pt);}}}
\newcommand{\mysetminusT}{\mysetminusD}
\newcommand{\mysetminusS}{\hbox{\tikz{\draw[line width=0.45pt,line cap=round] (2pt,0) -- (0,4pt);}}}
\newcommand{\mysetminusSS}{\hbox{\tikz{\draw[line width=0.4pt,line cap=round] (1.5pt,0) -- (0,3pt);}}}

\renewcommand{\setminus}{\mathbin{\mathchoice{\mysetminusD}{\mysetminusT}{\mysetminusS}{\mysetminusSS}}}
